\documentclass[aspectratio=169]{beamer}

% ======================
% Packages
% ======================
\usepackage[utf8]{inputenc}
\usepackage[T1]{fontenc}
\usepackage[brazil]{babel}
\usepackage{amsmath, amssymb}
\usepackage{microtype}
\usepackage{csquotes}
\usepackage{natbib}
\usepackage{tikz}

% ======================
% Colors
% ======================
\definecolor{mblack}{RGB}{20, 20, 20}
\definecolor{mblue}{RGB}{0, 102, 204}
\definecolor{morange}{RGB}{230, 100, 20}
\definecolor{mwhite}{RGB}{255, 255, 255}
% ======================
% Theme & Style
% ======================
\usetheme{default}
\usecolortheme{default}
\setbeamercolor{normal text}{fg=mblack, bg=white}
\setbeamercolor{frametitle}{fg=white, bg=mblue}
\setbeamercolor{title}{fg=white}
\setbeamercolor{author}{fg=white}
\setbeamercolor{date}{fg=white}
\setbeamercolor{itemize item}{fg=morange}
\setbeamercolor{itemize subitem}{fg=mblue}
\setbeamercolor{enumerate item}{fg=morange}
\setbeamercolor{enumerate subitem}{fg=mblue}
\setbeamercolor{block title}{fg=white, bg=mblue}
\setbeamercolor{block body}{fg=mblack, bg=mblue!8}
\setbeamercolor{alerted text}{fg=morange}

\setbeamerfont{frametitle}{series=\bfseries, size=\large}
\setbeamerfont{title}{series=\bfseries, size=\LARGE}
\setbeamerfont{block title}{series=\bfseries}

% Remove navigation symbols
\setbeamertemplate{navigation symbols}{}

% Footer
\setbeamertemplate{footline}{%
  \leavevmode%
  \hbox{%
    \begin{beamercolorbox}[wd=0.7\paperwidth, ht=2.25ex, dp=1ex, leftskip=0.3cm]{normal text}%
      \tiny\color{mblack!50} BCC502 -- Metodologia Científica para Ciência da Computação
    \end{beamercolorbox}%
    \begin{beamercolorbox}[wd=0.3\paperwidth, ht=2.25ex, dp=1ex, rightskip=0.3cm, right]{normal text}%
      \tiny\color{mblack!50}\insertframenumber{} / \inserttotalframenumber
    \end{beamercolorbox}%
  }%
}

% Frametitle with left orange accent bar
\setbeamertemplate{frametitle}{%
  \nointerlineskip%
  \begin{beamercolorbox}[wd=\paperwidth, ht=0.55cm, dp=0.2cm, leftskip=0.5cm]{frametitle}%
    \usebeamerfont{frametitle}\insertframetitle%
  \end{beamercolorbox}%
  \vspace{-0.1cm}%
  \textcolor{morange}{\rule{\paperwidth}{2pt}}%
}

% Title page colors
\setbeamercolor{titlebar top}{bg=mwhite}
\setbeamercolor{titlebar main}{bg=mblue}
\setbeamercolor{titlebar bottom}{bg=morange}

\setbeamertemplate{title page}{%
  \vspace{0pt}%
  \begin{beamercolorbox}[wd=\paperwidth, ht=0.55cm, dp=0pt]{titlebar top}%
  \end{beamercolorbox}%
  \begin{beamercolorbox}[wd=\paperwidth, ht=3.8cm, dp=0pt, center]{titlebar main}%
    \vspace{0.6cm}
    {\usebeamerfont{title}\color{white}\inserttitle\par}%
    \vspace{0.35cm}
    {\color{white}\small\insertauthor\par}%
    \vspace{0.15cm}
    {\color{white}\footnotesize\insertdate\par}%
    \vspace{0.25cm}
  \end{beamercolorbox}%
  \begin{beamercolorbox}[wd=\paperwidth, ht=1.3cm, dp=0pt]{titlebar bottom}%
  \end{beamercolorbox}%
}

% Item spacing
\setlength{\itemsep}{4pt}

% Custom commands
\newcommand{\impt}[1]{\textcolor{morange}{\textbf{#1}}}
\newcommand{\ntc}[1]{\textcolor{mblue}{#1}}

% ======================
% Document Info
% ======================
\title{BCC502 -- Metodologia Científica}
\author{Introdução ao Indutivismo \footnote{Baseado em \citep{chalmers1993ciencia}}}
\date{\today}

% ======================
% Document
% ======================
\begin{document}

% --- Title Slide ---
\begin{frame}[plain]
  \titlepage 
\end{frame}

% --- Table of Contents ---
\begin{frame}{Sumário}
  \tableofcontents
\end{frame}

% ==============================================================
\section{Introdução}
% ==============================================================

\begin{frame}{Introdução}
  \begin{itemize}
    \item O que é ciência?
    \pause
    \item O que queremos dizer quando usamos a palavra \textit{``científica''} como adjetivo?
    \pause
    \begin{itemize}
      \item Parece implicar algum tipo de mérito.
      \item Algum tipo especial de confiabilidade.
      \item Por quê?
      \item Existe um ``método científico'' que atribui à ciência estas características?
    \end{itemize}
    \pause
    \item Este é o trabalho principal de um filósofo da ciência.
    \pause
    \item Vamos estudar os principais resultados da filosofia da ciência.
    \pause
    \begin{itemize}
      \item Perspectiva histórica.
      \item Começamos com as ideias mais rudimentares e vamos expandindo para as ideias mais complexas e modernas.
    \end{itemize}
  \end{itemize}
\end{frame}

% ==============================================================
\section{Ciência: Concepção de Senso Comum}
% ==============================================================

\begin{frame}{Ciência: Concepção de \textit{Senso Comum}}
  \begin{itemize}
    \item Conhecimento científico é conhecimento \textbf{provado}.
    \item As teorias científicas são derivadas de maneira rigorosa dados adquiridos por observação e experimentação.
    \item A ciência é baseada no que podemos ver, ouvir, tocar, etc.
    \item Opiniões ou preferências pessoais e suposições especulativas \textbf{não têm lugar} na ciência.
    \item A ciência é \textbf{objetiva}.
    \item O conhecimento científico é confiável porque é conhecimento provado objetivamente.
  \end{itemize}
\end{frame}

\begin{frame}{Ciência: Concepção de \textit{Senso Comum}}
  Esta concepção se tornou popular no século XVII como consequência da revolução científica da época (Galileu, Newton).

  \vspace{0.4cm}
  \begin{itemize}
    \item ``Se quisermos compreender a natureza, devemos consultar a natureza''.
    \item É um erro utilizar obras dos antigos (Aristóteles, Bíblia, \ldots) como fonte de conhecimento.
    \item A experiência é a fonte do conhecimento.
    \item A ciência é uma estrutura construída sobre \textbf{fatos}.
  \end{itemize}
\end{frame}

\begin{frame}{Galileu e a Ruptura com a Tradição}
  \begin{block}{Citação}
    ``Não foram tanto as observações e experimentos de Galileu que causaram a ruptura com a tradição, mas sua atitude em relação a eles. Para ele, os dados eram tratados como dados, e não relacionados a alguma ideia preconcebida\ldots{} Os dados da observação poderiam ou não se adequar a um esquema conhecido do universo, mas a coisa mais importante, na opinião de Galileu, era aceitar os dados e construir a teoria para se adequar a eles.''
  \end{block}
  \vspace{0.2cm}
  \hfill\citep{anthony1957science}
\end{frame}

% ==============================================================
\section{Indutivismo}
% ==============================================================

\begin{frame}{Indutivismo}
  \begin{center}
    \large\ntc{Ciência como conhecimento derivado dos dados da experiência}
  \end{center}
\end{frame}

\begin{frame}{Indutivista Ingênuo}
  A Ciência Começa com a Observação
  \begin{itemize}
    \item O observador científico deve ter órgãos sensitivos normais e inalterados e registra fielmente o que pode ver, ouvir, etc.
    \pause
    \item \textbf{Sem preconceitos.}
    \pause
    \item Afirmações (proposições) a respeito do estado do mundo podem ser justificadas de maneira direta pelo uso dos sentidos do observador não-preconceituoso.
    \pause
    \item Proposições de observação formam a base a partir da qual leis e teorias científicas devem ser derivadas.
  \end{itemize}
\end{frame}

\begin{frame}{Proposições de Observação}
  \begin{itemize}
    \item \textbf{Exemplos:} (i) Hoje temos uma lua cheia. (ii) O sol nasceu no leste.
    \pause
    \item \textbf{Características}:
    \begin{itemize}
      \item A verdade de tais observações deve ser estabelecida com cuidadosa observação.
      \item Qualquer observador pode estabelecer ou conferir sua verdade pelo uso direto de seus sentidos.
      \item Observadores podem ver por si mesmos. 
    \end{itemize}
  \end{itemize}
\end{frame}

\begin{frame}{Tipos de Afirmação}
  \begin{itemize}
    \item \textbf{\ntc{Afirmações singulares}}: Referem-se a uma ocorrência específica ou a um estado de coisas num lugar específico, num tempo específico. \ntc{Proposições de observação são afirmações singulares.}
    \pause
    \vspace{0.3cm}
    \item \textbf{\ntc{Afirmações universais}}: São informações gerais que afirmam coisas sobre as propriedades ou comportamento de algum aspecto do universo.
    \begin{itemize}
      \item Referem-se a todos os eventos de um tipo específico em todos os lugares em todos os tempos. Ex.: \textit{O Sol nasce no leste.}
      \item \impt{As leis que constituem o conhecimento científico são sempre afirmações universais.}
    \end{itemize}
  \end{itemize}
  \pause
  \vspace{0.3cm}
  \begin{alertblock}{Questão central}
    Se a ciência é baseada na experiência, então por que meios é possível extrair das afirmações singulares, que resultam da observação, as afirmações universais, que constituem o conhecimento científico?
  \end{alertblock}
\end{frame}

\begin{frame}{Resposta Indutivista}
  É legítimo generalizar afirmações universais a partir de observações e afirmações singulares desde que as condições a seguir sejam satisfeitas:

  \vspace{0.3cm}
  \begin{enumerate}
    \item O número de proposições de observação (afirmações singulares) que forma a base da generalização deve ser \textbf{grande};
    \item As observações devem ser repetidas sob \textbf{ampla variedade de condições};
    \item \textbf{Nenhuma} proposição de observação deve conflitar com a afirmação universal derivada.
  \end{enumerate}
\end{frame}

\begin{frame}{O Princípio Indutivista}
  \begin{block}{Princípio}
    Se um grande número de $A$s foi observado sob uma ampla variedade de condições, e se todos os $A$s observados possuíam sem exceção a propriedade $B$, então todos os $A$s têm a propriedade $B$.
  \end{block}
\end{frame}

% ==============================================================
\section{Explicações e Previsões}
% ==============================================================

\begin{frame}{Explicações e Previsões}
  Não basta para a ciência gerar afirmações universais --- é fundamental que a partir das afirmações universais se gerem \textbf{explicações} e \textbf{previsões}.

  \vspace{0.4cm}
  O tipo de raciocínio envolvido neste tipo de derivação é o \ntc{raciocínio dedutivo}.
\end{frame}

\begin{frame}{Raciocínio Dedutivo}
  \begin{block}{Definição}
    Um argumento dedutivo é \textit{logicamente válido} se, supondo que as premissas são verdadeiras, a conclusão é necessariamente verdadeira --- isto é, é impossível que as premissas sejam verdadeiras e a conclusão falsa.
  \end{block}

  \vspace{0.3cm}
  \begin{alertblock}{Atenção}
    Validade \textbf{não} garante que as premissas sejam de fato verdadeiras. Um argumento válido com premissas verdadeiras é chamado de \textit{sólido} (\textit{sound}).
  \end{alertblock}
\end{frame}

\begin{frame}{Raciocínio Dedutivo: Exemplo 1}
  \textbf{Premissas:}
  \begin{enumerate}
    \item Água razoavelmente pura congela a cerca de $0^\circ$C.
    \item O radiador do meu carro contém água razoavelmente pura.
  \end{enumerate}

  \vspace{0.4cm}
  \textbf{Conclusão:}
  \begin{itemize}
    \item[\(\Rightarrow\)] Se a temperatura cair abaixo de $0^\circ$C, a água do radiador do meu carro vai congelar.
  \end{itemize}
\end{frame}

\begin{frame}{Raciocínio Dedutivo: Exemplo 2}
  \textbf{Premissas:}
  \begin{enumerate}
    \item Todos os gatos têm 5 patas.
    \item Oswaldo é meu gato.
  \end{enumerate}

  \vspace{0.4cm}
  \textbf{Conclusão:}
  \begin{itemize}
    \item[\(\Rightarrow\)] Oswaldo tem 5 patas.
  \end{itemize}

  \vspace{0.3cm}
  \textit{\small Argumento válido, mas \impt{não sólido} --- a premissa 1 é falsa.}
\end{frame}

\begin{frame}{Padrão do Raciocínio Dedutivo}
  \begin{columns}
    \begin{column}{0.48\textwidth}
      \textbf{Premissas:}
      \begin{enumerate}
        \item Leis e teorias
        \item Condições iniciais
      \end{enumerate}
    \end{column}
    \begin{column}{0.48\textwidth}
      \textbf{Conclusão:}
      \begin{itemize}
        \item[\(\Rightarrow\)] Previsões e explicações
      \end{itemize}
    \end{column}
  \end{columns}
  \pause
  \vspace{0.5cm}
  \begin{alertblock}{Síntese do Indutivismo}
    No indutivismo, leis, teorias e condições iniciais são obtidas por observação direta e indução. Então previsões e explicações podem ser deduzidas deles.
  \end{alertblock}
\end{frame}

% ==============================================================
\section{A Atração do Indutivismo Ingênuo}
% ==============================================================

\begin{frame}{A Atração do Indutivismo Ingênuo}
  \begin{enumerate}
    \item A maioria das pessoas vai achar a explicação indutivista da ciência adequada.
    \item Ela fornece uma explicação formalizada de algumas das impressões popularmente mantidas a respeito da ciência:
    \begin{itemize}
      \item Poder de explicação e previsão
      \item Objetividade
      \item Confiabilidade superior
    \end{itemize}
  \end{enumerate}
\end{frame}

\begin{frame}{Poder de Explicar e Prever}
  \begin{itemize}
    \item Leis e teorias podem ser obtidas por \textbf{indução} a partir de afirmações singulares (proposições de observação).
    \item Condições iniciais são proposições de afirmação e podem ser obtidas por \textbf{observação direta}.
    \item Explicações e previsões podem ser obtidas por \textbf{dedução}.
  \end{itemize}
\end{frame}

\begin{frame}{Objetividade}
  \begin{itemize}
    \item Observação e raciocínio indutivos são processos \textbf{objetivos}.
    \item Proposições de observação podem ser averiguadas por qualquer observador pelo uso normal dos sentidos. Assim, não podem depender do gosto, da opinião ou das esperanças do observador.
    \item As induções satisfazem ou não as condições prescritas --- não é uma questão \textbf{subjetiva}.
  \end{itemize}
\end{frame}

\begin{frame}{Confiabilidade}
  \begin{enumerate}
    \item Seguem das definições indutivistas sobre observação e indução.
    \item As proposições de afirmação formam a base da ciência e são \textbf{seguras e confiáveis}.
    \item Essa confiabilidade é \textbf{transmitida} às leis e teorias derivadas, desde que as condições de indução sejam satisfeitas.
  \end{enumerate}
\end{frame}

% --- References ---
\begin{frame}[allowframebreaks]{Referências}
  \bibliographystyle{apalike}
  \bibliography{references}
\end{frame}

\end{document}
